\makeatletter
\def\input@path{{../.format/}}
\makeatother

\documentclass[runningheads]{llncs}
\usepackage[T1]{fontenc}
\usepackage{graphicx}

\begin{document}

% Anonymous proxy for \maketitle
\begin{center}
	{\Large\bfseries\boldmath
		On the Formal Inexplicability of Self-Evident Metaphysical Phenomena and Related Systems%
		\par}
	\vskip 0.8cm
\end{center}

\begin{abstract}
Since Descartes first proclaimed "cogito, ergo sum" in 1637, "I think, therefore I am" has become the Declaration of Independence for consciousness. We hold these truths to be self-evident that our subjective experiences are real and require no further justification. This philosophical fact is fortunate for us because consciousness appears to resist scientific explanation of any form. Chalmers' codified this observation into a Hard Problem that is essentially \textit{Hard} by identifying the explanatory gap that reveals a hidden incompleteness in science: no physical account can fully explain how neural processes give rise to the subjective quality of conscious experience. We are forced to confront a strange paradox: the most certain aspect of our existence eludes our most powerful explanatory framework. Recently, a group referring to themselves as the \textit{IIT-Concerned} clarified their stance that a theory of consciousness that cannot be tested will simply be deemed "unscientific," underscoring the tension between the current state of affairs in the field and the demands of a framework akin to actual science. The emergence of Semantic AI has transformed what was once a philosophical oddity into a critical practical concern, with profound ethical implications. We simply cannot assess whether AI systems are conscious without first understanding the necessary and sufficient conditions that give rise to consciousness. How can we provide a formal definition of conscious experience that preserves the essence of its metaphysical nature in a framework that is physically falsifiable? 
		
\keywords{Formal Theories of Consciousness \and The Hard Problem \and Subjectivity \and Semantic AI \and Metaphysical Transduction.}
\end{abstract}

%-----------------------------------------------------------------------------------------------------

\begin{thebibliography}{99}

\bibitem{Descartes1637}
Descartes, R.: \textit{Discourse on Method}. (1637)

\bibitem{Chalmers1995}
Chalmers, D.J.: Facing Up to the Problem of Consciousness. \textit{Journal of Consciousness Studies} \textbf{2}(3), 200--219 (1995)

\bibitem{IITConcerned2023}
IIT-Concerned et al.: What makes a theory of consciousness unscientific? \textit{Nature Neuroscience} (2023)  
\url{https://doi.org/10.1038/s41593-023-01439-9}

\end{thebibliography}
\end{document}
